\newpage
\datos{color}{}{}
\section*{Definición multi-objetivo del problema y resolución}
\addcontentsline{toc}{section}{Definición multi-objetivo del problema y resolución}

\subsection*{Fundamento técnico y justificación}
Hasta este punto, el problema se ha abordado mediante una \textbf{función de mérito penalizada} que 
combina objetivos heterogéneos (tiempo máximo $T_{max}$, desincronización $|t_A - t_B|$ y 
violaciones de restricciones) en un único valor escalar. Según las recomendaciones 
técnicas recibidas, este enfoque presenta debilidades críticas: una alta sensibilidad a los 
pesos asignados y la posibilidad de generar soluciones artificiales por compensación numérica 
entre términos de distinta naturaleza física. \\

Por tanto, como se adelantó en la propuesta inicial, el problema del \textit{rendezvous} robótico 
definido en este trabajo se puede formular como \textbf{bi-objetivo}. El objetivo ideal es 
minimizar simultáneamente los tiempos de ejecución de ambos robots:
\begin{equation}
\min_{\mathbf{p}_A,\mathbf{q}_A,\mathbf{p}_B,\mathbf{q}_B,\mathbf{s}_A,\mathbf{s}_B,\mathbf{s}_{rail}}
\;\; \mathbf{T}(\cdot),
\end{equation}

con
\[
\mathbf{T}(\cdot) = 
\begin{bmatrix}
t_A(\mathbf{p}_A,\mathbf{q}_A,\mathbf{s}_A) \\
t_B(\mathbf{p}_B,\mathbf{q}_B,\mathbf{s}_B,\mathbf{s}_{rail})
\end{bmatrix},
\]

%\begin{equation}
%    \min \mathbf{f}(\mathbf{x}) = (t_A, t_B)
%\end{equation}
sujeto a la restricción de sincronización $|t_A - t_B| \le \Delta t$ y a las restricciones 
geométricas y cinemáticas del sistema. El planteamiento multiobjetivo permite obtener 
un \textbf{frente de Pareto}, proporcionando un conjunto de soluciones óptimas que revelan el 
compromiso real (\textit{trade-off}) entre la velocidad de cada robot y la precisión de la 
sincronización, en lugar de una única solución condicionada por un ajuste manual de pesos.

\subsection*{Algoritmo seleccionado: NSGA-II}
Para la resolución de este modelo se ha seleccionado el algoritmo \textbf{NSGA-II} 
(\textit{Nondominated Sorting Genetic Algorithm II}). Esta técnica es el estándar para problemas 
con múltiples objetivos debido a su eficiencia y robustez. El algoritmo se basa en dos mecanismos 
fundamentales:
\begin{itemize}
    \item \textbf{Ordenación No Dominada (\textit{Nondominated Sorting}):} Clasifica a los individuos en 
    frentes sucesivos de dominancia de Pareto. Una solución domina a otra si es igual de buena en 
    todos los objetivos y estrictamente mejor en al menos uno.
    \item \textbf{Distancia de Apiñamiento (\textit{Crowding Distance}):} Se utiliza para desempatar entre 
    soluciones del mismo frente, favoreciendo aquellas situadas en regiones menos densas del espacio
    de objetivos para garantizar una distribución uniforme de soluciones a lo largo de todo el frente
    de Pareto.
\end{itemize}

\subsection*{Diseño y operadores}
Siguiendo la arquitectura de dos niveles validada en el Tema 2, se mantiene una
\textbf{codificación real} de las variables de decisión: $\mathbf{x} = (s_A, s_B, s_{rail}, p_{rail})$. 
Los operadores se adaptan de la siguiente manera:
\begin{itemize}
    \item \textbf{Cruce y Mutación:} Se emplean operadores para codificación real compatibles con el
    sistema de cobots, manteniendo la reparación de soluciones (\textit{InRange}) tras cada 
    operación para asegurar la factibilidad cinemática.
    \item \textbf{Selección de Padres:} Se utiliza el \textbf{método de torneo}, que resulta
    más sencillo y eficaz en problemas multiobjetivo al permitir comparaciones directas basadas 
    en la definición de optimalidad de Pareto.
    \item \textbf{Recombinación Elitista:} NSGA-II combina la población actual con la descendencia 
    generada y selecciona los mejores individuos para la siguiente generación basándose en el 
    frente de Pareto y la diversidad, asegurando que los mejores hallazgos no se pierdan.
\end{itemize}

\subsection*{Plan experimental y métricas de calidad}
Siguiendo el rigor estadístico del \textbf{Anexo de evaluación}, se realizarán \textbf{30 ejecuciones independientes} para capturar la variabilidad del proceso \cite{27, 198}. Dado que el resultado es un conjunto de soluciones (frente de Pareto) y no un valor único, se emplearán métricas específicas de optimización multiobjetivo \cite{180, 199}:
\begin{enumerate}
    \item \textbf{Hipervolumen (HV):} Mide el volumen del espacio de objetivos dominado por el frente obtenido respecto a un punto de referencia \cite{180, 184}.
    \item \textbf{Indicador Epsilon ($\epsilon$):} Evalúa la proximidad del frente hallado respecto a un frente de referencia ideal \cite{180, 183}.
    \item \textbf{Diversidad y Cobertura:} Grado de extensión y uniformidad de las soluciones a lo largo de la frontera de Pareto \cite{180, 199}.
\end{enumerate}

\subsection*{Resultados (Esquema)}

Los resultados obtenidos mediante la aplicación de NSGA-II se presentan a continuación, comparando la calidad del frente de Pareto obtenido frente a los resultados mono-objetivo del algoritmo de Evolución Diferencial previo.

\begin{table}[h!]
\centering
\begin{tabular}{|l|c|c|c|c|}
\hline
\textbf{Configuración} & \textbf{Media HV} & \textbf{Desv. Típica HV} & \textbf{Mejor $\epsilon$} & \textbf{Factibilidad (\%)} \\ \hline
NSGA-II (Propuesto)    & --                & --                       & --                        & --                         \\ \hline
DE Mono-objetivo       & --                & --                       & --                        & --                         \\ \hline
\end{tabular}
\caption{Métricas de rendimiento multiobjetivo comparadas.}
\end{table}

\begin{figure}[h!]
    \centering
    % Aquí se incluiría el gráfico del Frente de Pareto (t_A vs t_B)
    \caption{Visualización del Frente de Pareto obtenido por NSGA-II.}
\end{figure}
